% id: cameron-martin
% title: An Application of the Cameron-Martin Theorem
% date: 15 May. 2024

\documentclass{article}
\usepackage{amsmath, amssymb}

\begin{document}

The Cameron-Martin theorem provides a characterization of the Gaussian measure on infinite dimensional
spaces. While I first saw this definition in Martin's lectures on SPDEs, it was mostly presented
with context with its application to SPDEs. In this entry, I'll record an application of this
theorem in the context of obtaining a maximal inequality for Brownian motions.

The Cameron-Martin theorem is the following.

\begin{theorem}
Let \(\mu \in \mathcal{M}(X)\) be a Gaussian and denote \(\mathcal{H}_\mu\) for its Cameron-Martin space.
For \(h \in X\), denoting \(\tau_h : x \in X \mapsto x + h\), we have that \(\tau_h^* \mu \ll \mu\) if and only
if \(h \in \mathcal{H}_\mu\). In this case,
\[\frac{d\tau_h^* \mu}{d\mu}(x) = \exp\left(h^*(x) - \frac{1}{2}\|h\|^2_\mu\right).\]
\end{theorem}

Let \(\mu\) be the Gaussian measure on \(\mathcal{C}([0, 1], \mathbb{R})\) with covariance
\(C_\mu(\delta_t, \delta_s) = s \wedge t\). \(\mu\) corresponds to the law of the Brownian motion
on \([0, 1]\) and we call it the Wiener measure. We have that \(\mathcal{H}_\mu = H^{1, 2}_0([0, 1])\) -
the space of all absolutely continuous functions \(h\) with \(h(0) = 0\) and \(\dot h(s) \in L^2([0, 1])\).
Moreover, for \(h \in \mathcal{H}_\mu\),
\(h^*(k) = \int_0^1 \dot h(s) \dot k(s) d s\). As this defines a bounded linear functional from \(\mathcal{H}_\mu\),
we can almost surely uniquely extend \(h^*\) to a measurable element of \(C_\mu(\delta_t, \delta_s)^*\).
To see this, it suffices to check that
\[h_t = \delta_t(h) = C_\mu(\delta_t, h^*) = h^*(C_\mu(\delta_t)).\]
We see that
\[C_\mu(\delta_t)_s = \left(\int \omega \delta_t(\omega) \mu(d \omega)\right)_s
= \int \omega_s \omega_t \mu(d \omega) = C_\mu(\delta_t, \delta_s) = s \wedge t.\]
Thus, \(\dot{C}_\mu(\delta_t)_s = \mathbb{1}_{[0, t]}(s)\) and hence,
\[h^*(C_\mu(\delta_t)) = \int_0^t \dot f_s d s = f_t\]
as required.

As an application, the Cameron-Martin theorem allows us to bound the probability that a Brownian motion
stays near a given function. In particular, for \(f \in H^{1, 2}_0([0, 1])\), denoting \(\mu_W\) for the
Wiener measure, we have that
\[\mathbb{P}(|B_t - f|_\infty \le \epsilon) = \tau_{-f}^* \mu_W(\overline B_\epsilon(0))
= e^{-\frac{1}{2}\|f\|^2_{H^{1, 2}_0([0, 1])}} \int_{\overline B_\epsilon(0)} e^{- f^*(\omega)} \mu_W(d \omega).\]
This is then bounded above by
\[e^{-\frac{1}{2}\|f\|^2_{H^{1, 2}_0([0, 1])} + \epsilon \|f^*\|_{X^*}} \mathbb{P}(|B_t|_\infty \le \epsilon).\]
On the other hand, using the inequality \(e^x \ge 1 + x\), we have that
\[\int_{\overline B_\epsilon(0)} e^{- f^*(\omega)} \mu_W(d \omega) \ge
\int_{\overline B_\epsilon(0)} (1 + - f^*(\omega)) \mu_W(d \omega)
\ge \mathbb{P}(|B_t|_\infty \le \epsilon) - \int_{\overline B_\epsilon(0)} f^*(\omega) \mu_W(d \omega)\]
Now, since \(f^*_* \mu_W \sim \mathcal{N}(0, \|f\|_{\mathcal{H}_W}^2)\), the last integral is zero by spherical symmetry.
Consequently, have the bound
\[\mathbb{P}(|B_t - f|_\infty \le \epsilon) \in
e^{-\frac{1}{2}\|f\|^2_{H^{1, 2}_0([0, 1])}}
\left[\mathbb{P}(|B_t|_\infty \le \epsilon), e^{\epsilon \|f^*\|_{X^*}}\mathbb{P}(|B_t|_\infty \le \epsilon)\right]\]
and thus,
\[\lim_{\epsilon \downarrow 0} \frac{\mathbb{P}(|B_t - f|_\infty \le \epsilon)}{\mathbb{P}(|B_t|_\infty \le \epsilon)}
= e^{-\frac{1}{2}\|f\|^2_{H^{1, 2}_0([0, 1])}}.\]

\end{document}
